% !TEX root = ssh-indep.tex
% ssh-indep.tex - version autonome du TP SSH (test PR)
% Preambule inline (copie de ../../inc/encommun.tex) : aucune dependance externe,
% le fichier se compile seul (latexmk -pdf -shell-escape ssh-indep.tex), n'importe ou.
% doc sur la class exam
% http://fr.wikibooks.org/wiki/LaTeX/LaTeX_pour_enseignants
% http://agregationchimie.free.fr/exam.php


%QUESTIONS ET REPONSES

\documentclass[a4paper,10pt]{exam}
\usepackage[utf8]{inputenc}
\usepackage[T1]{fontenc}
\usepackage[french]{babel}


\frenchsetup{AutoSpacePunctuation=false}
\usepackage{graphicx}

\author{Pouchoulon}

%recodage numerotation \part
\renewcommand{\partlabel}{\thepartno}

\newcommand{\iconetext}[1]{
\begin{array}{c}
\includegraphics{#1}
\end{array}}

\newcommand{\jmp}{Jean-Marc Pouchoulon}
\graphicspath{{../../images/}}

%coloration de code

\usepackage[table,xcdraw]{xcolor}
\usepackage{listings}


\definecolor{mygreen}{rgb}{0,0.6,0}
\definecolor{mygray}{rgb}{0.5,0.5,0.5}
\definecolor{mymauve}{rgb}{0.58,0,0.82}
\definecolor{mygreen}{rgb}{0,0.6,0}
\definecolor{mygrey}{rgb}{0.5,0.5,0.5}
\definecolor{mymauve}{rgb}{0.58,0,0.82}
\definecolor{carolinablue}{rgb}{0.6, 0.73, 0.89}
\definecolor{vertpomme}{rgb}{0.66, 0.89, 0.63}
\definecolor{jaune}{rgb}{0.91, 0.84, 0.42}
\definecolor{cerise}{rgb}{0.96, 0.0, 0.63}
\definecolor{bg}{rgb}{0.85,0.85,0.85}
\lstset{ %
  backgroundcolor=\color{white},   % choose the background color; you must add \usepackage{color} or \usepackage{xcolor}
  basicstyle=\footnotesize,        % the size of the fonts that are used for the code
  breakatwhitespace=false,         % sets if automatic breaks should only happen at whitespace
  breaklines=true,                 % sets automatic line breaking
  captionpos=b,                    % sets the caption-position to bottom
  commentstyle=\color{mygreen},    % comment style
  deletekeywords={...},            % if you want to delete keywords from the given language
  escapeinside={\%*}{*)},          % if you want to add LaTeX within your code
  extendedchars=true,              % lets you use non-ASCII characters; for 8-bits encodings only, does not work with UTF-8
  frame=single,                    % adds a frame around the code
  keepspaces=true,                 % keeps spaces in text, useful for keeping indentation of code (possibly needs columns=flexible)
  keywordstyle=\color{blue},       % keyword style
  language=Octave,                 % the language of the code
  morekeywords={*,...},            % if you want to add more keywords to the set
  numbers=left,                    % where to put the line-numbers; possible values are (none, left, right)
  numbersep=5pt,                   % how far the line-numbers are from the code
  numberstyle=\tiny\color{mygray}, % the style that is used for the line-numbers
  rulecolor=\color{black},         % if not set, the frame-color may be changed on line-breaks within not-black text (e.g. comments (green here))
  showspaces=false,                % show spaces everywhere adding particular underscores; it overrides 'showstringspaces'
  showstringspaces=false,          % underline spaces within strings only
  showtabs=false,                  % show tabs within strings adding particular underscores
  stepnumber=2,                    % the step between two line-numbers. If it's 1, each line will be numbered
  stringstyle=\color{mymauve},     % string literal style
  tabsize=2,                       % sets default tabsize to 2 spaces
  title=\lstname                   % show the filename of files included with \lstinputlisting; also try caption instead of title
  }

% eviter les boites rouges dans pygmentize quand le parser détecte un problème

%\usemintedstyle{xcode}
%\usemintedstyle{solarized-dark}
%\usemintedstyle{colorful}
%\usemintedstyle{autumn}
%\usepackage{fancyhdr}


% Éviter les carrés rouges
\usepackage[hyperfootnotes=false]{hyperref}
\usepackage{minted2}
\newminted{docker}{
    linenos=false,
    bgcolor=bg,
    baselinestretch=1.2, 
    frame=none,
    framesep=10pt,
    %    tabsize=4,
    %  gobble=8,
    fontfamily=helvetica,
    fontsize=\footnotesize,
}
\newminted{ldif}{
    linenos=false,
    bgcolor=bg,
    baselinestretch=1.2, 
    frame=none,
    framesep=10pt,
    %    tabsize=4,
    %  gobble=8,
    fontfamily=helvetica,
    fontsize=\footnotesize,
}

\newminted{yaml}{
    linenos=false,
    bgcolor=bg,
    baselinestretch=1.2, 
    frame=none,
    framesep=10pt,
    showtabs=false,
    %    tabsize=4,
    %  gobble=8,
    fontfamily=helvetica,
    fontsize=\footnotesize,
}

\newminted{bash}{
    linenos=false,
    bgcolor=bg,
    baselinestretch=1.2, 
    frame=none,
    framesep=10pt,
    %    tabsize=4,
    %  gobble=8,
    fontfamily=helvetica,
    fontsize=\footnotesize,
}


\newminted{basht}{
    linenos=false,
    bgcolor=bg,
    baselinestretch=1.2, 
    frame=none,
    framesep=10pt,
    %    tabsize=4,
    %  gobble=8,
    fontfamily=helvetica,
    fontsize=\tiny,
}
\newminted{bashg}{
    linenos=false,
    bgcolor=bg,
    baselinestretch=1.2, 
    frame=none,
    framesep=10pt,
    %    tabsize=4,
    %  gobble=8,
    fontfamily=helvetica,
    fontsize=\footnotesize,
}

\newminted{json}{
    linenos=false,
    bgcolor=bg,
    baselinestretch=1.2, 
    frame=none,
    framesep=10pt,
    %    tabsize=4,
    %  gobble=8,
    fontfamily=helvetica,
    fontsize=\tiny,
}

\newminted{text}{
    linenos=false,
    baselinestretch=1.2, 
    %    bgcolor=bg,
    %    frame=none,
    %    framesep=10pt,
    %    tabsize=4,
    %  gobble=8,
    fontfamily=helvetica,
   % fontsize=\tiny,
}


\newminted{powershell}{
    breaklines=true, 
    linenos=false,
    bgcolor=bg,
    baselinestretch=1.2, 
    frame=none,
    framesep=10pt,
    %    tabsize=4,
    %  gobble=8,
    fontfamily=helvetica,
    fontsize=\footnotesize,
}


\newminted{powershellsession}{
    linenos=false,
    bgcolor=bg,
    baselinestretch=1.2, 
    frame=none,
    framesep=10pt,
    %    tabsize=4,
    %  gobble=8,
    fontfamily=helvetica,
    fontsize=\footnotesize,
}
\newminted{batch}{
    linenos=false,
    breaklines=true, 
    bgcolor=bg,
    baselinestretch=1.2, 
    frame=none,
    framesep=10pt,
    %    tabsize=4,
    %  gobble=8,
    fontfamily=helvetica,
    fontsize=\footnotesize,
}
\newminted{ini}{
    linenos=false,
    breaklines=true, 
    bgcolor=bg,
    baselinestretch=1.2, 
    frame=none,
    framesep=10pt,
    %    tabsize=4,
    %  gobble=8,
    fontfamily=helvetica,
    fontsize=\footnotesize,
}
\newminted{python}{	
    breaklines=true,
    linenos=true,
    bgcolor=bg,
    baselinestretch=1.2, 
    frame=none,
    framesep=10pt,
    %    tabsize=4,
    %  gobble=8,
    fontfamily=helvetica,
    fontsize=\footnotesize,
}
\newminted{ansible-output}{	
    breaklines=true,
    linenos=true,
    bgcolor=bg,
    baselinestretch=1.2, 
    frame=none,
    framesep=10pt,
    %    tabsize=4,
    %  gobble=8,
    fontfamily=helvetica,
    fontsize=\footnotesize,
}
\newminted{jinja}{	
    breaklines=true,
    linenos=true,
    bgcolor=bg,
    baselinestretch=1.2, 
    frame=none,
    framesep=10pt,
    %    tabsize=4,
    %  gobble=8,
    fontfamily=helvetica,
    fontsize=\footnotesize,
}
\newminted{toml}{
	linenos=false,
	baselinestretch=1.2, 
	%    bgcolor=bg,
	%    frame=none,
	%    framesep=10pt,
	%    tabsize=4,
	%  gobble=8,
	fontfamily=helvetica,
%	fontsize=\tiny,
}

\newminted{js}{
    linenos=false,
    bgcolor=white,
    bgcolor=bg,
    baselinestretch=1.2, 
    frame=none,
    framesep=10pt,
%    tabsize=4,
%  gobble=8,
    fontfamily=helvetica,
    fontsize=\footnotesize,
}
\newminted{terraform}{
    linenos=false,
    bgcolor=white,
    bgcolor=bg,
    baselinestretch=1.2, 
    frame=none,
    framesep=10pt,
%    tabsize=4,
%  gobble=8,
    fontfamily=helvetica,
    fontsize=\footnotesize,
}
\newminted{ios}{
    linenos=false,
    bgcolor=white,
    bgcolor=bg,
    baselinestretch=1.2, 
    frame=none,
    framesep=10pt,
%    tabsize=4,
%  gobble=8,
    fontfamily=helvetica,
    fontsize=\footnotesize,
}
\newminted{text2}{
    linenos=false,
    baselinestretch=1.2, 
%    bgcolor=bg,
%    frame=none,
%    framesep=10pt,
%    tabsize=4,
%  gobble=8,
    fontfamily=helvetica,
%    fontsize=\tiny,
}
%\usemintedstyle{murphy}
\usepackage{enumitem}

\usemintedstyle{tango}
\author{Jean-Marc Pouchoulon}
%\author{\IEEEauthorblockN{Jean-Marc Pouchoulon}
%\IEEEauthorblockA{Département Réseaux et Télécoms\\
%IUT Béziers\\
%Université de Montpellier\\
%Email: jean-marc.pouchoulon@iutbeziers.fr}}

% pour afficher ( ou pas) les réponses dans le document
%\printanswers

\usepackage{parskip}

\usepackage{microtype}
\DisableLigatures{encoding = *, family = * } % Désactive toutes les ligatures
% no red boxes on parser error:
\makeatletter
\expandafter\def\csname PYGdefault@tok@err\endcsname{\def\PYGdefault@bc##1{{\strut ##1}}}
\makeatother

\title{TP R202 : SSH --- TIPS \& TRICKS\\{\small clés, agent, bastion, tunnels, multiplexage}}

\date{Octobre 2026}
%%\author{Jean-Marc Pouchoulon}

% pour afficher (ou pas) les réponses dans le document
\printanswers
\noprintanswers

\graphicspath{{./images/}}


\begin{document}
\maketitle

\section{Objectifs}

\textbf{SSH} (Secure SHell) est le protocole de communication chiffrée le plus utilisé pour
administrer des serveurs à distance. Au-delà du simple \texttt{ssh serveur}, il regorge de
fonctionnalités qui, une fois maîtrisées, changent radicalement l'efficacité d'un administrateur :
authentification par clés, tunnels, connexions par saut, multiplexage.

À la fin de ce TP, vous devez être capable de :
\begin{itemize}
  \item générer, déployer, révoquer des clés \textbf{ed25519} et comprendre les permissions exigées ;
  \item utiliser un \textbf{agent SSH} (Linux et Windows) et expliquer le piège du \texttt{ForwardAgent} ;
  \item écrire un \texttt{\textasciitilde/.ssh/config} qui supprime toute option tapée en ligne de commande ;
  \item se connecter à un serveur interne \textbf{via un bastion} avec \texttt{ProxyJump} ;
  \item ouvrir des tunnels \texttt{-L}, \texttt{-R} et \texttt{-D} et savoir précisément \emph{qui écoute où} ;
  \item activer le \textbf{multiplexage} \texttt{ControlMaster} et mesurer son impact ;
  \item déboguer une connexion : \texttt{-v}, \texttt{known\_hosts}, \texttt{StrictHostKeyChecking}.
\end{itemize}
Vous aurer besoin de deux VMs : un bastion (exposé sur le port 2222) 
et une cible (accessible uniquement depuis le bastion).

\section{Clés : bien faire les choses dès le départ}

L'authentification par clé repose sur un pair asymétrique : la clé privée reste sur votre poste,
la clé publique est déposée dans \texttt{\textasciitilde/.ssh/authorized\_keys} du serveur.
Depuis OpenSSH 6.5 (2014), l'algorithme recommandé est \textbf{Ed25519} (RFC 8032) :
courbes elliptiques, clé de 256~bits, aussi sûre qu'une RSA 3072, plus courte, plus rapide,
et sans paramètre à « bien choisir » comme la taille d'une RSA.

\begin{questions}
\question Générez sur votre poste une paire de clés ed25519 avec le commentaire
\texttt{"prenom.nom@poste"} et une phrase de passe. Quels fichiers sont créés ?
Quelles permissions portent-ils, et pourquoi cette différence ?
\begin{solution}
\begin{bashcode}
cd ~/.ssh
ssh-keygen -t ed25519 -C "prenom.nom@poste"
# -f permet de choisir le nom de fichier (sinon id_ed25519)
# -N "" supprime la phrase de passe (à éviter !)
ls -l id_ed25519 id_ed25519.pub
# -rw------- 1 ... id_ed25519      (600 : privée, lisible par son propriétaire seul)
# -rw-r--r-- 1 ... id_ed25519.pub  (644 : publique, partageable librement)
\end{bashcode}
Le \texttt{.pub} n'a \emph{rien} de secret : c'est lui que l'on dépose sur les serveurs.
La clé privée est chiffrée avec votre phrase de passe ; une permission 644 sur la clé privée
provoque un refus immédiat de la part du client (voir question suivante).
\end{solution}
\question Affichez l'empreinte (fingerprint) de la clé avec \texttt{ssh-keygen -lf}.
Que représente cette empreinte ? Pourquoi est-elle utile lors d'un audit d'un
\texttt{authorized\_keys} sur un serveur ?
\begin{solution}
\begin{bashcode}
ssh-keygen -lf id_ed25519.pub
# 256 SHA256:abcdEFGH... prenom.nom@poste (ED25519)
ssh-keygen -lf id_ed25519     # marche aussi sur la clé privée
\end{bashcode}
C'est le SHA-256 de la clé publique : court, stable, unique. Sur un serveur où
\texttt{authorized\_keys} compte des dizaines de clés opaques, le fingerprint permet
d'identifier \emph{qui} a accès sans avoir à recoller les clés deux par deux, et de
cross-checker la clé réellement déployée avec celle de l'émetteur (empreinte annoncée
par un canal séparé).
\end{solution}
\question Que fait \texttt{ssh-keygen -y -f id\_ed25519} ? En quoi est-ce une astuce
de vérification ?
\begin{solution}
\begin{bashcode}
ssh-keygen -y -f id_ed25519        # redérive la publique depuis la privée
ssh-keygen -y -f id_ed25519 | diff - <(cut -d' ' -f1-2 id_ed25519.pub) && echo CONFORME
\end{bashcode}
On peut donc vérifier, côté client, qu'une ligne d'un \texttt{authorized\_keys}
loin (récupérée par \texttt{cat}) correspond bien à \emph{votre} clé privée, sans
confiance dans l'affichage de la clé publique seule.
\end{solution}
\question Déployez la clé sur le bastion. D'abord avec l'outil dédié, puis décrivez la
méthode manuelle. Quelles permissions exactes \texttt{sshd} exige-t-il sur
\texttt{\textasciitilde}, \texttt{\textasciitilde/.ssh} et \texttt{authorized\_keys}
(option \texttt{StrictModes}, activée par défaut) ?
\begin{solution}
\begin{bashcode}
# depuis l'hôte (mot de passe la dernière fois) :
ssh-copy-id -p 2222 root@127.0.0.1
# méthode manuelle :
cat ~/.ssh/id_ed25519.pub | ssh -p 2222 root@127.0.0.1 \
  'mkdir -p ~/.ssh && cat >> ~/.ssh/authorized_keys && chmod 700 ~/.ssh && chmod 600 ~/.ssh/authorized_keys'
ssh -p 2222 root@127.0.0.1   # plus de mot de passe demandé
\end{bashcode}
Avec \texttt{StrictModes yes} (défaut), \texttt{sshd} refuse l'authentification par clé si
le \texttt{home}, \texttt{~/.ssh} ou \texttt{authorized\_keys} sont \emph{group/world}
écrivables ou non possédés par l'utilisateur (700/600 attendus ; 644 toléré sur
\texttt{authorized\_keys}, 600 recommandé).
\end{solution}
\question Piège : faites \texttt{chmod 644} sur votre clé privée, puis essayez de vous connecter.
Que répond le client ? Pourquoi le \emph{client} (et non le serveur) refuse-t-il ?
\begin{solution}
\begin{bashcode}
chmod 644 ~/.ssh/id_ed25519
ssh -p 2222 root@127.0.0.1
# @@@@@@@@@@@@@@@@@@@@@@@@@@@@@@@@@@@@@@@@@@@@@@@@@@@@@@@@@@@
# @         UNPROTECTED PRIVATE KEY FILE!          @
# ...
# Bad permissions. Private key should be readable
#     only by you.
# Permission denied (publickey).
chmod 600 ~/.ssh/id_ed25519   # la remise en état
\end{bashcode}
Le client applique ce contrôle \emph{avant même de charger} la clé : la clé privée est
chiffrée avec la phrase de passe, mais le client ne veut pas \emph{demander} la phrase de
passe si d'autres utilisateurs de la machine peuvent lire le fichier (la phrase de passe
saisie en clair sur le terminal serait alors utilisée pour décrypter une clé lisible par
n'importe qui). C'est une défense en profondeur côté client, indépendante du serveur.
\end{solution}
\question Générez une \emph{seconde} clé, \texttt{id\_ed25519\_cible}, qui ne sera déposée
\textbf{que} sur la cible. Quel option permet de préciser quelle clé présenter ?
Pourquoi la « une clé par contexte » est-elle une bonne hygiène ?
\begin{solution}
\begin{bashcode}
ssh-keygen -t ed25519 -f ~/.ssh/id_ed25519_cible -C "prenom.nom@poste-pour-cible"
ssh -i ~/.ssh/id_ed25519_cible ...    # sur la ligne de commande
IdentityFile ~/.ssh/id_ed25519_cible  # dans le config (section suivante)
\end{bashcode}
Avantages : la révocation est chirurgicale (on retire une ligne d'un
\texttt{authorized\_keys} sans impacter les autres accès) ; l'audit de journal montre
quelle clé a été utilisée ; une clé perdue/compromis ne donne accès qu'à un périmètre ;
on peut même utiliser des algorithmes ou durées de vie différents selon le contexte.
\end{solution}
\question Un stagiaire a quitté le service : sa clé est encore dans
\texttt{authorized\_keys} du bastion. Révoquez-la et montrez dans les journaux du bastion
comment identifier quelle clé un utilisateur vient d'utiliser pour se connecter.
\begin{solution}
\begin{bashcode}
# sur le bastion :
nano ~/.ssh/authorized_keys     # retirer la ligne correspondante (vérifier le fingerprint avant !)
# puis vérifier dans les journaux :
journalctl -u ssh | grep -i "accepted"
# Sep 29 10:12:03 bastion sshd[1234]: Accepted publickey for root
#   (ED25519 SHA256:abcdEFGH...) from 10.99.0.1 port 52344 ssh2:
#   RSA SHA256:...             <- la clé du client
#   ED25519 SHA256:abcdEFGH... <- la clé acceptée, à croiser avec ssh-keygen -lf
\end{bashcode}
La révocation d'une clé SSH se fait \emph{à la main} (pas de CRL) : retirer la ligne de
\texttt{authorized\_keys} est l'opération. D'où l'intérêt des commentaires lisibles et
des fingerprints.
\end{solution}
\question Comparez : générez une RSA 3072 et mesurez la taille des fichiers par rapport à
votre ed25519. En déduisez un argument taille/sécurité à transmettre à un collègue
qui « fait du RSA 4096 pour être plus sûr ».
\begin{solution}
\begin{bashcode}
ssh-keygen -t rsa -b 3072 -f ~/.ssh/id_rsa_test -N "" -q
ls -l id_ed25519 id_rsa_test
# id_ed25519  ~400 octets   id_rsa_test ~2600 octets
\end{bashcode}
Ed25519 offre un niveau de sécurité équivalent à RSA 3072 pour une clé 6 fois plus
petite, une signature plus rapide, et sans choix de paramètre (pas de « 2048 suffit ?
non, 3072 ? 4096 ? »). RSA 4096 n'est « plus sûr » que ed25519 que de très peu, au prix
de 10 fois le poids.
\end{solution}
\question Que fait \texttt{ssh-keyscan} ? À quoi sert-il dans l'automatisation, et quelle
est la limite \textbf{sécurité} fondamentale à connaître avant de l'utiliser ?
\begin{solution}
\begin{bashcode}
# pré-peupler known_hosts d'une flotte (ex. avant un run d'automatisation) :
ssh-keyscan -t ed25519 -p 2222 127.0.0.1 >> ~/.ssh/known_hosts
ssh-keyscan -f serveurs.txt          # une IP/hostname par ligne
\end{bashcode}
Il récupère les clés \emph{publiques de hôte} sans s'authentifier. Limite fondamentale :
il ne fournit \textbf{aucune authentification} du serveur — sur le premier contact, un
MITM peut répondre avec \emph{sa} clé, qui sera alors enregistrée comme bonne.
À n'utiliser qu'après vérification de l'empreinte par un canal de confiance
(out-of-band), ou en environnement maîtrisé.
\end{solution}
\end{questions}

\section{L'agent : la phrase de passe, une seule fois}

\texttt{ssh-agent} est un démon qui conserve vos clés privées \emph{déjà déverrouillées} en
mémoire. Une fois la phrase de passe saisie, toutes les connexions (et les scripts qui
lancent \texttt{ssh}/\texttt{scp}/\texttt{rsync}/\texttt{git} en local) se reconnectent sans
y repenser.

\begin{questions}
\question Que fait exactement \texttt{ssh-add} ? Quelles variables d'environnement relient
votre session à l'agent ?
\begin{solution}
\begin{bashcode}
echo $SSH_AUTH_SOCK $SSH_AGENT_PID
ssh-add                    # demande la phrase de passe, charge la clé
ssh -p 2222 root@127.0.0.1 # plus de phrase de passe
\end{bashcode}
\texttt{SSH\_AUTH\_SOCK} pointe vers le socket UNIX de l'agent, \texttt{SSH\_AGENT\_PID}
vers son PID. Tout processus de la session (votre shell, un éditeur, un script) qui voit
cette variable peut demander à l'agent de signer — c'est ce qui rend
\texttt{git clone git@github.com:} ou \texttt{rsync} transparents.
\end{solution}
\question Chargez vos deux clés dans l'agent. Listez-les.
\begin{solution}
\begin{bashcode}
ssh-add ~/.ssh/id_ed25519 ~/.ssh/id_ed25519_cible
ssh-add -l
# 256 SHA256:aaaa... prenom.nom@poste (ED25519)
# 256 SHA256:bbbb... prenom.nom@poste-pour-cible (ED25519)
\end{bashcode}
\end{solution}
\question Limitez la durée de vie d'une clé à 1~heure. Comment observez-vous son expiration ?
\begin{solution}
\begin{bashcode}
ssh-add -t 3600 ~/.ssh/id_ed25519_cible
ssh-add -l          # présente pendant 1h
# ... une heure plus tard ...
ssh-add -l          # la clé a disparu de la liste, il faudra re-ssh-add
\end{bashcode}
\texttt{-t} prend une durée en secondes. Utile sur un poste partagé ou en salle de travaux
pratiques : la clé quitte la mémoire de l'agent toute seule.
\end{solution}
\question Si votre poste de travail est \textbf{Windows} (client OpenSSH natif) : le service
\texttt{ssh-agent} joue le même rôle. Quelles commandes PowerShell vérifient le service et
chargent une clé ?
\begin{solution}
\begin{bashcode}
Get-Service ssh-agent
Start-Service ssh-agent          # si arrêté
ssh-add $HOME\.ssh\id_ed25519
ssh-add -l
\end{bashcode}
Compte-rendu utile à retenir : sur Windows, les clés des comptes \emph{administrateurs}
peuvent aussi être placées dans
\texttt{C:\textbackslash ProgramData\textbackslash ssh\textbackslash administrators\_authorized\_keys}
côté serveur, mais côté client rien ne change : c'est le même \texttt{ssh-add}.
\end{solution}
\question Ajoutez \texttt{AddKeysToAgent yes} dans votre config (section suivante) et expliquez
le comportement obtenu.
\begin{solution}
Dès qu'une clé déverrouillée \emph{a servi} pour s'authentifier, elle est automatiquement
ajoutée à l'agent pour la suite de la session : vous ne re-saisissez la phrase de passe
qu'une fois, même si vous aviez oublié de faire \texttt{ssh-add} avant.
\end{solution}
\question \textbf{Le piège du \texttt{ForwardAgent}.} Connectez-vous au bastion avec
\texttt{ssh -A} (ou \texttt{ForwardAgent yes}), puis exécutez \emph{sur le bastion}
\texttt{ssh-add -l}. Que voyez-vous ? Formulez précisément le risque, puis expliquez
pourquoi le \texttt{ProxyJump} (section suivante) ne présente pas ce risque.
\begin{solution}
\begin{bashcode}
ssh -A -p 2222 root@127.0.0.1
# sur le bastion :
ssh-add -l
# les clés de VOTRE agent local sont listées sur le bastion !
\end{bashcode}
Risque : pendant la durée de votre session, \emph{tout} processus ou utilisateur ayant un
accès shell sur le bastion peut utiliser \textbf{vos} clés (la variable \texttt{SSH\_AUTH\_SOCK}
est transmise dans la session). Si le bastion est compromis, vos accès à \emph{tous} les
serveurs où ces clés sont déposées le sont aussi.
Avec \texttt{ProxyJump}, la clé est présentée \emph{depuis votre poste} pour chaque saut :
elle ne transite jamais \emph{sur} le bastion ; le bastion n'a qu'un simple flux TCP
transité. Principe : chaque saut a sa propre clé, rien ne se transite.
\end{solution}
\end{questions}

\section{Le fichier \texttt{\textasciitilde/.ssh/config}}

La ligne de commande \texttt{ssh} est faite pour être tapée ; le fichier de configuration est
fait pour qu'on ne tape \emph{plus rien}. Règle essentielle : la première valeur trouvée
pour une option est celle qui est utilisée, donc les hôtes spécifiques doivent précéder les
blocs \texttt{Host *}.

\begin{questions}
\question Écrivez un \texttt{\textasciitilde/.ssh/config} définissant deux alias,
\texttt{bastion} et \texttt{cible} (avec \texttt{ProxyJump bastion} pour la cible),
chacun avec son \texttt{IdentityFile}. Vérifiez que \texttt{ssh bastion} fonctionne.
\begin{solution}
\begin{bashcode}
cat > ~/.ssh/config <<'EOF'
Host bastion
    HostName 127.0.0.1
    Port 2222
    User root
    IdentityFile ~/.ssh/id_ed25519

Host cible
    HostName 10.99.0.11
    User root
    IdentityFile ~/.ssh/id_ed25519_cible
    ProxyJump bastion
EOF
chmod 600 ~/.ssh/config
ssh bastion && ssh cible
\end{bashcode}
Désormais : plus d'IP, plus de \texttt{-p}, plus de \texttt{-i}, plus de
\texttt{-o ProxyCommand=...}.
\end{solution}
\question Complétez avec un bloc \texttt{Host *} de valeurs par défaut
(\texttt{ServerAliveInterval}, \texttt{ConnectTimeout}, \texttt{AddKeysToAgent}).
Quelles options \textbf{ne doivent jamais} figurer dans un \texttt{Host *} et pourquoi ?
\begin{solution}
\begin{bashcode}
cat > ~/.ssh/config <<'EOF'
Host *
    ServerAliveInterval 60        # keepalive client->serveur : détecte les connexions mortes
    ServerAliveCountMax 3
    ConnectTimeout 10
    AddKeysToAgent yes
    ControlMaster auto
    ControlPath ~/.ssh/ctl-%C
    ControlPersist 10m

Host bastion
    HostName 127.0.0.1
    Port 2222
    User root
    IdentityFile ~/.ssh/id_ed25519

Host cible
    HostName 10.99.0.11
    User root
    IdentityFile ~/.ssh/id_ed25519_cible
    ProxyJump bastion
EOF
\end{bashcode}
Jamais dans \texttt{Host *} : \texttt{ForwardAgent} (le piège de la section précédente
s'appliquerait à \emph{tous} vos serveurs) et \texttt{IdentityFile} global (présenter
systématiquement la même clé partout revient à ne pas isoler les contextes, et rend la
révocation par clé impossible à cibler).
\end{solution}
\question Si un hôte liste \emph{plusieurs} \texttt{IdentityFile}, quel est le comportement
d'OpenSSH ? Et que change \texttt{IdentitiesOnly yes} ?
\begin{solution}
Sans \texttt{IdentitiesOnly}, OpenSSH essaie \emph{toutes} les clés candidates
(IdentityFiles de la config \emph{plus} toutes les clés de l'agent) jusqu'à ce qu'une
soit acceptée. C'est pratique, mais ça expose votre agent au serveur : chaque clé
présentée est annoncée. \texttt{IdentitiesOnly yes} restreint la tentative aux
\texttt{IdentityFile} explicites uniquement — utile quand un serveur est saturé de
tentatives ou quand on veut un comportement déterministe (c'est aussi l'un des pièges
du challenge final).
\end{solution}
\question Astuce de diagnostic : \texttt{ssh -G} affiche la configuration \emph{résolue}
pour un hôte, sans se connecter. Utilisez-le pour vérifier quel \texttt{IdentityFile}
s'applique à \texttt{cible} et quelle valeur effective a \texttt{ServerAliveInterval}.
\begin{solution}
\begin{bashcode}
ssh -G cible | grep -Ei "^(identityfile|serveraliveinterval|proxyjump) "
# identityfile /home/pouchou/.ssh/id_ed25519_cible
# serveraliveinterval 60
# proxyjump bastion
\end{bashcode}
\texttt{ssh -G} est l'outil de « debug de config » par excellence : il matérialise la
fusion Host */Host spécifique et les valeurs par défaut.
\end{solution}
\question Votre parc grandit : vous voulez partager un fichier de config d'équipe sans
écraser le vôtre. Utilisez la directive \texttt{Include} (OpenSSH $\geq$ 7.4) pour charger
\texttt{~/.ssh/config.d/labs.conf}.
\begin{solution}
\begin{bashcode}
mkdir -p ~/.ssh/config.d
cat > ~/.ssh/config.d/labs.conf <<'EOF'
Host lab-*
    User root
    StrictHostKeyChecking accept-new
EOF
# en tête du config principal :
sed -i '1i Include config.d/*.conf' ~/.ssh/config
ssh -G lab-42 | grep -i user
\end{bashcode}
\texttt{Include} est évalué à l'endroit où il est écrit (donc en tête = haute priorité).
C'est le mécanisme standard pour les dotfiles : le config « personnel » reste versionné
séparément des configs « contexte ».
\end{solution}
\end{questions}

\section{Bastion : la connexion par saut (\texttt{ProxyJump})}

Scénario type d'entreprise : le serveur \texttt{cible} n'écoute que sur le réseau interne
; le \textbf{bastion} est la seule machine exposée. Objectif : \texttt{ssh cible} en une
commande, avec la bonne clé à chaque saut, \emph{sans jamais transmettre de clé au bastion}.

\begin{questions}
\question Précisez la différence \textbf{sémantique} entre \texttt{ProxyJump} et
\texttt{ForwardAgent}.
\begin{solution}
\texttt{ForwardAgent} : l'agent \emph{local} est \emph{disponible} sur le saut (socket
transféré) ; c'est le client distant qui choisit et présente les clés — la clé, elle,
n'est pas copiée, mais elle est \emph{utilisable} depuis le saut.
\texttt{ProxyJump} : le client \emph{local} établit lui-même les deux connexions SSH
(saut puis cible) ; le saut ne voit qu'un flux TCP brut qu'il transite (\texttt{-W}).
Aucune clé ne passe sur le saut, aucune variable d'agent n'y est exposée.
\end{solution}
\question \textbf{Ancienne syntaxe} : connectez-vous à la cible avec
\texttt{ssh -o "ProxyCommand=ssh bastion -W \%h:\%p" 10.99.0.11}. Qui exécute quoi sur
le bastion ? Que fait exactement \texttt{-W hôte:port} ?
\begin{solution}
\begin{bashcode}
ssh -o "ProxyCommand=ssh bastion -W %h:%p" 10.99.0.11
# %h -> 10.99.0.11, %p -> 22
\end{bashcode}
La \texttt{ProxyCommand} s'exécute \emph{en local} ; ici elle lance un \texttt{ssh} vers
le bastion en mode \texttt{-W} : le daemon du bastion ouvre un flux TCP vers
\texttt{\%h:\%p} et y branche le stdin/stdout du client local. Le bastion ne crée
\textbf{aucune session shell} pour la connexion à la cible — il transite.
\end{solution}
\question \textbf{Syntaxe moderne} : faites la même chose avec \texttt{-J}. Pourquoi est-ce
plus sûr et plus lisible ?
\begin{solution}
\begin{bashcode}
ssh -J bastion 10.99.0.11
# ou avec l'alias :
ssh -J bastion cible
\end{bashcode}
\texttt{-J} est un \texttt{ProxyCommand} standardisé : OpenSSH construit lui-même la
chaîne, gère les erreurs saut/cible séparément, s'intègre au config, et compose avec
\texttt{-L}/\texttt{-R} (les tunnels s'appliquent sur la destination finale). En ligne de
commande, c'est aussi infiniment plus court et auditable.
\end{solution}
\question Comment OpenSSH choisit-il la clé \textbf{pour le saut} et la clé \textbf{pour
la destination} ? (Indice : votre agent contient les deux clés.)
\begin{solution}
Pour chaque connexion de la chaîne (saut, puis cible), OpenSSH tente les clés candidates
de l'agent et les \texttt{IdentityFile} de la config \emph{de cet hôte} :
\begin{itemize}
  \item le saut \texttt{bastion} : la config de \texttt{bastion} pointe sur
        \texttt{id\_ed25519} $\rightarrow$ c'est la clé du bastion qui est présentée ;
  \item la destination \texttt{cible} : \texttt{IdentityFile id\_ed25519\_cible} $\rightarrow$
        la clé cible.
\end{itemize}
Comme la clé \texttt{cible} est absente du bastion et la clé \texttt{bastion} absente de la
cible, même un contournement par l'agent ne permet pas d'aller plus loin que le périmètre
de chaque clé. \emph{Chaque saut a sa clé} : c'est le modèle de sécurité à retenir.
\end{solution}
\question Sans modifier la configuration, sur une ligne de commande unique exécutée
depuis votre hôte, obtenez \texttt{hostname \&\& whoami} sur la cible via le bastion.
\begin{solution}
\begin{bashcode}
ssh -J bastion 10.99.0.11 'hostname && whoami'
# cible
# root
\end{bashcode}
\end{solution}
\question Preuve de transit : depuis la cible, listez les connexions TCP établies
(\texttt{ss -tn state established | grep 22}). Qui est l'adresse source des sessions
\texttt{ssh} ? Que voit \texttt{sshd} du bastion quand vous vous connectez à la cible ?
\begin{solution}
\begin{bashcode}
# sur la cible :
ss -tn state established | grep :22
# TCP  ESTAB  10.99.0.11:22    10.99.0.10:39412     <- la source est le BASTION
\end{bashcode}
Pour \texttt{sshd} de la cible, la connexion vient de \texttt{10.99.0.10} : le bastion.
C'est le modèle classique de journalisation d'un bastion (le bastion doit conserver ses
propres journaux pour rattacher « qui, depuis l'extérieur » à « quelle connexion interne »).
\end{solution}
\end{questions}

\section{Tunnels : du forwarding à la demande}

Trois mécanismes, un même principe : des flux TCP \emph{embarqués} dans la connexion SSH
chiffrée. Rappels de topologie : l'hôte ne voit que le bastion (2222) ; le service 8080
tourne sur la cible.

\begin{questions}
\question \textbf{Tunnel local} : depuis l'hôte, ouvrez
\texttt{ssh -N -L 18080:127.0.0.1:8080 cible} et récupérez le service avec
\texttt{curl 127.0.0.1:18080}. Décomposez la syntaxe \texttt{-L
[adresse\_locale:]port\_local:hote:port} et répondez à la question capitale :
\textbf{sur quelle machine} \texttt{hote:port} est-elle interprétée ?
\begin{solution}
\begin{bashcode}
ssh -N -L 18080:127.0.0.1:8080 cible
curl 127.0.0.1:18080
# SERVICE INTERNE CIBLE 10.99.0.11
\end{bashcode}
\texttt{127.0.0.1:8080} est résolu \textbf{du point de vue de la destination} de la
connexion SSH (la cible, ou plus généralement le dernier hôte de la chaîne si on passe
par un \texttt{ProxyJump}) : le daemon SSH de la cible ouvre la connexion locale vers son
\texttt{127.0.0.1:8080}. C'est pour cette raison qu'on écrit
\texttt{-L 5432:localhost:5432} pour une base : « la base qui tourne \emph{sur le serveur
où je me connecte} ».
\texttt{-N} = pas de session shell distante, juste le tunnel.
\end{solution}
\question Qui \emph{écoute} sur le port 18080 de votre hôte ? Montrez-le avec
\texttt{ss -ltnp}. Que se passe-t-il si vous ajoutez
\texttt{-L 0.0.0.0:18080:127.0.0.1:8080} ? Et quelle est la politique par défaut d'OpenSSH
récente ?
\begin{solution}
\begin{bashcode}
ss -ltnp | grep 18080
# LISTEN  0  128  127.0.0.1:18080  0.0.0.0:*  users:(("ssh",pid=...))
\end{bashcode}
C'est le \textbf{processus client \texttt{ssh} de votre poste} qui écoute, et qui reboucle
sur la cible. Avec l'adresse explicite \texttt{0.0.0.0}, l'écoute devient visible sur
\emph{toutes} les interfaces de l'hôte : n'importe qui sur le LAN de l'hôte peut atteindre
le service interne via vous — à ne faire que pour un usage contrôlé (ex. : partage vers
une autre machine locale). Depuis OpenSSH 7.4, l'écoute locale est par défaut
\textbf{restreinte à la boucle} : c'est le comportement sûr.
\end{solution}
\question Cas professionnel : accéder à une base PostgreSQL interne \emph{sans ouvrir le
port 5432}. Écrivez la commande (tunnel local sur 5433) et l'outil client qui l'utilise.
\begin{solution}
\begin{bashcode}
ssh -N -f -L 5433:127.0.0.1:5432 cible
# -f : passe en arrière-plan après l'authentification (nécessite -N ou -L/-R/-D)
pg_dump -h 127.0.0.1 -p 5433 -U postgres ma_base > dump.sql
psql -h 127.0.0.1 -p 5433 -U postgres
\end{bashcode}
Avantages sur l'ouverture directe du port : le service n'a pas à écouter sur une
interface externe (souvent \texttt{listen\_addresses = localhost} dans
\texttt{postgresql.conf}), pas de règle de firewall supplémentaire, le chiffrement est
celui de SSH, et la révocation se fait en coupant le tunnel. C'est \emph{la} méthode
standard pour les bases, l'UI de gestion, l'API interne.
\end{solution}
\question \textbf{Tunnel distant} : sur votre \emph{hôte}, lancez
\texttt{python3 -m http.server 3000} dans un répertoire contenant un fichier
\texttt{marqueur.html}. Depuis l'hôte, ouvrez
\texttt{ssh -R 8080:127.0.0.1:3000 cible}, puis depuis \emph{la cible}
\texttt{curl 127.0.0.1:8080}. Qui écoute où cette fois ?
\begin{solution}
\begin{bashcode}
# sur l'hôte :
mkdir -p /tmp/www && echo "MARQUEUR DEPUIS L'HOTE" > /tmp/www/marqueur.html
cd /tmp/www && python3 -m http.server 3000 &
ssh -R 8080:127.0.0.1:3000 cible
# sur la cible :
curl 127.0.0.1:8080/marqueur.html
# MARQUEUR DEPUIS L'HOTE
\end{bashcode}
Symétrique du local : c'est le \textbf{daemon \texttt{sshd} de la cible} qui écoute sur
\texttt{127.0.0.1:8080} et renvoie le flux vers \texttt{127.0.0.1:3000} \emph{de l'hôte}.
Usages typiques : exposer à distance un service qui tourne sur votre poste (Jupyter, un
dashboard, un outil de debug), ou recevoir le trafic d'un serveur qui ne peut pas sortir.
\end{solution}
\question Côté serveur, quelles options \texttt{sshd\_config} gouvernent ces tunnels
(\texttt{AllowTcpForwarding}, \texttt{GatewayPorts}) ? Que bloque
\texttt{AllowTcpForwarding no} sur le bastion, et pourquoi un « bastion kiosque » peut
être configuré ainsi ?
\begin{solution}
\begin{itemize}
  \item \texttt{AllowTcpForwarding yes} (défaut) : autorise les tunnels \emph{locaux}
        (\texttt{-L}) depuis les clients.
  \item \texttt{AllowTcpForwarding no} : refuse les tunnels \texttt{-R} \emph{demandés
        par le client distant} (le client local \texttt{-L} est quand même possible ;
        c'est \texttt{PermitTunnel} et les règles \texttt{Match} qui affinent).
        Sur un bastion « kiosque » (accès shell only, pas de tunnel),
        \texttt{AllowTcpForwarding no} empêche d'utiliser le bastion comme relais pour
        contourner le réseau.
  \item \texttt{GatewayPorts no} (défaut) : les écouteurs \texttt{-R} ne se bindent que sur
        la boucle de la machine distante. \texttt{GatewayPorts clients/all} autoriserait
        le bind sur d'autres interfaces : à réserver à des cas très encadrés.
\end{itemize}
\end{solution}
\question \textbf{Tunnel dynamique} : ouvrez \texttt{ssh -N -D 1080 cible} puis depuis
l'hôte : \texttt{curl --socks5-hostname 127.0.0.1:1080
http://10.99.0.11:8080/}. Expliquez ce que fait \texttt{-D}, la différence entre
\texttt{--socks5} et \texttt{--socks5-hostname}, et ce que le dynamique \emph{ne remplace
pas}.
\begin{solution}
\begin{bashcode}
ssh -N -D 1080 cible
curl --socks5-hostname 127.0.0.1:1080 http://10.99.0.11:8080/
# SERVICE INTERNE CIBLE 10.99.0.11
\end{bashcode}
\texttt{-D} ouvre un \textbf{proxy SOCKS5} : n'importe quel client SOCKS peut, à travers
le tunnel, se connecter à \emph{n'importe quelle} destination joignable \emph{depuis le
serveur} — un « VPN de fortune » au niveau applicatif.
\texttt{--socks5} résout le DNS \emph{en local} puis envoie l'IP ;
\texttt{--socks5-hostname} envoie le nom d'hôte et laisse le serveur résoudre : c'est la
bonne option car c'est la \emph{résolution côté serveur} qui compte (le DNS interne).
Limites à connaître : pas de multicast, pas d'accès \emph{entrant} sur l'hôte, pas de
transparence pour les applications non-SOCKS, et le chiffrement ne protège que le tunnel
SSH (le HTTPS reste chiffré de bout en bout). Pour du réseau complet, un vrai VPN
(WireGuard) reste l'outil adéquat.
\end{solution}
\question Astuce de durcissement : vous voulez interdire à l'équipe de créer des tunnels
depuis le bastion pendant une période (audit). Quelle option \texttt{sshd\_config} et quel
\texttt{Match} combiner ? (Pas besoin de tester, écrivez le bloc.)
\begin{solution}
\begin{bashcode}
# /etc/ssh/sshd_config (bastion)
Match Address 10.99.0.0/24
    AllowTcpForwarding no
# puis :
systemctl reload ssh
\end{bashcode}
\texttt{Match} applique les directives suivantes \emph{seulement} aux connexions qui y
correspondent (adresse source, utilisateur, hôte, clé...) — c'est l'outil de
différenciation fine du daemon.
\end{solution}
\end{questions}

\section{Multiplexage : \texttt{ControlMaster}, l'art du jour}

Chaque \texttt{ssh} classique coûte : une poignée de main TCP, un échange de clés
(chiffrement), une authentification complète. Avec le \textbf{multiplexage}, la première
connexion établit un canal maître ; toutes les suivantes réutilisent ce canal chiffré et
authentifié : les sessions ultérieures sont de simples \emph{channels} supplémentaires.

Votre config contient déjà (section \texttt{Host *} de la question 2) :
\begin{bashcode}
ControlMaster auto
ControlPath ~/.ssh/ctl-%C
ControlPersist 10m
\end{bashcode}

\begin{questions}
\question Mesurez l'effet : exécutez \textbf{deux fois de suite}
\texttt{time ssh bastion true}. Comparez les durées. Expliquez \emph{exactement} ce qui
est économisé à la deuxième exécution.
\begin{solution}
\begin{bashcode}
time ssh bastion true
# première : 0.3 à 2 s (TCP + kex + auth)
time ssh bastion true
# seconde : ~20-50 ms (ouverture d'un channel sur la connexion existante)
\end{bashcode}
Économisé à la seconde exécution : la TCP SYN/ACK, la négociation des algorithmes
(kex), l'échange des clés de session, et toute l'authentification (échange de
signatures). Il ne reste que l'ouverture d'un channel multiplexé sur la connexion
maître déjà chiffrée.
\end{solution}
\question Inspectez le maître : \texttt{ssh -O check bastion}. Que renseigne-t-il ?
Quels sont les autres \texttt{-O} utiles ?
\begin{solution}
\begin{bashcode}
ssh -O check bastion
# Master running (pid=12345), 1 channel opened.
#  Open master channel: 0/1 (persistent)
ssh -O cancel bastion     # tue le maître et ses tunnels
\end{bashcode}
\texttt{check} (état), \texttt{cancel} (arrêt), \texttt{logout} (déconnexion du maître
sans tuer les channels ouverts). Indispensable pour le dépannage « mon tunnel ne marche
plus, est-ce que le maître est encore vivant ? ».
\end{solution}
\question Observez le socket : \texttt{ls -l ~/.ssh/ctl-*}. Quel type de fichier, quelles
permissions ? Pourquoi sont-elles essentielles à la sécurité du multiplexage ?
\begin{solution}
\begin{bashcode}
ls -l ~/.ssh/ctl-*
# srw------- 1 pouchou pouchou 0 Sep 29 10:15 /home/pouchou/.ssh/ctl-abc123...
\end{bashcode}
C'est un \textbf{socket UNIX} (le \texttt{s}), en \texttt{0600} dans un
\texttt{~/.ssh} en \texttt{0700}. Quiconque peut \emph{écrire} sur ce socket peut
\emph{s'asseyer} sur la connexion maître : ouvrir des channels, exécuter des commandes,
ouvrir des tunnels \textbf{sans se re-authentifier}. D'où l'exigence : propriétaire seul,
dans un répertoire privé. Ne pointez jamais \texttt{ControlPath} vers
\texttt{/tmp} (monde lisible/écrivable).
\end{solution}
\question Pourquoi \texttt{\%C} plutôt que \texttt{\%r@\%h:\%p} dans \texttt{ControlPath} ?
\begin{solution}
\texttt{\%C} est un \emph{hash} de la combinaison \texttt{\%r@\%h:\%p} : la longueur du
chemin est donc bornée. Sur Linux, le chemin d'un socket UNIX est limité à
\textbf{108~caractères} (champ \texttt{sun\_path}) ; avec de longs
\texttt{HostName} (FQDN) et de longs chemins de home,
\texttt{~/.ssh/ctl-\%r@\%h:\%p} déborde et la connexion échoue avec un message du type
\texttt{Socket path too long}. \texttt{\%C} élimine ce genre d'échec bête.
\end{solution}
\question \textbf{Tunnels qui survivent} : ouvrez un tunnel en arrière-plan
(\texttt{ssh -fN -L 18080:127.0.0.1:8080 cible}), puis \emph{fermez le terminal}.
Le tunnel est-il encore là ? Pourquoi le combo \texttt{ControlPersist} + \texttt{-fN} est-il
le bon réflexe ?
\begin{solution}
\begin{bashcode}
ssh -fN -L 18080:127.0.0.1:8080 cible
# ... fermeture du terminal ...
curl 127.0.0.1:18080     # toujours servable
ssh -O check cible       # Master running (pid=...)
\end{bashcode}
Sans \texttt{ControlPersist}, le maître meurt quand le dernier client se déconnecte : le
tunnel \texttt{-fN} mourrait avec votre terminal. Avec
\texttt{ControlPersist 10m}, le maître survit 10~minutes après la dernière déconnexion
(idem pour un tunnel \texttt{-fN}) : il est là quand vous revenez, et il se nettoie tout
seul si vous oubliez. C'est le combo « tunnel persistant mais pas zombie ».
\end{solution}
\question \texttt{ControlPersist} accepte une durée, \texttt{yes}, ou \texttt{unset}.
Que fait chacun ?
\begin{solution}
\begin{itemize}
  \item \texttt{ControlPersist 10m} : le maître survit 10~min après la dernière
        déconnexion, puis s'éteint ;
  \item \texttt{ControlPersist yes} : le maître persiste \emph{indéfiniment} (jusqu'à
        \texttt{-O cancel} ou au reboot) — pratique, mais à ne pas laisser par défaut ;
  \item \texttt{ControlPersist unset} : désactive la persistance ; le maître meurt avec
        le dernier client.
\end{itemize}
\end{solution}
\question Multiplexage et \textbf{scripts} : chronométrez un script qui ouvre 20 connexions
séquentielles vers le bastion avec et sans multiplexage. Citez un outil d'automatisation
qui profite de ce mécanisme.
\begin{solution}
\begin{bashcode}
time for i in $(seq 20); do ssh bastion true; done
# sans ControlMaster : ~30-60 s (20 auths complètes)
# avec            : ~0.5-1 s  (1 auth + 19 channels)
\end{bashcode}
Les outilleurs d'automatisation sur SSH (les playbooks qui ouvrent une session par tache,
certaines implémentations d'outils de provisionnement type pyinfra/ansible en mode
\texttt{ssh}) réduisent drastiquement leur temps de run quand le client multiplexe ;
c'est aussi le mécanisme que \texttt{git} (via libssh2/OpenSSH) et \texttt{rsync}
\texttt{ssh://} bénéficient indirectement.
\end{solution}
\question \textbf{Dépannage} : deux pannes classiques du multiplexage.
(a)~\texttt{Master has died} au milieu d'une session ; (b)~un maître « zombie » qui
n'accepte plus rien (réseau mort). Dans chaque cas, la commande à exécuter.
\begin{solution}
\begin{bashcode}
# (a) le serveur a coupé (reboot, keepalive dépassé) :
ssh -O check bastion     # "Error: Control socket ... does not exist"
ssh bastion              # une simple reconnexion rétablit un nouveau maître
# (b) un maître coincé (WiFi coupé, NAT expiré) :
ssh -O check bastion     # "Master running" mais les commandes partent en timeout
ssh -O cancel bastion    # purge le zombie ; la prochaine connexion en crée un frais
\end{bashcode}
Le réflexe de dépannage est toujours : \texttt{-O check} pour voir l'état,
\texttt{-O cancel} pour forcer, re-connexion.
\end{solution}
\end{questions}

\section{Exécuter à distance, déboguer, survivre}

\begin{questions}
\question Exécutez une commande sur la cible \emph{sans ouvrir de shell interactif} et
récupérez son code de retour localement. Montrez aussi le code de retour « spécial »
qu'OpenSSH utilise pour ses propres erreurs.
\begin{solution}
\begin{bashcode}
ssh cible 'df -h /'
ssh cible 'exit 42'; echo $?
# 42  -> le code de retour distant est propagé tel quel
ssh hote-inexistant true; echo $?
# 255 -> code réservé aux erreurs du client SSH lui-même
\end{bashcode}
\texttt{255} est la convention OpenSSH pour « c'est \emph{moi} le client qui a un
problème » (résolution, config, réseau), pas la commande distante.
\end{solution}
\question \textbf{Le piège du stdin} : le script ci-dessous ne s'exécute qu'une fois au
lieu de deux. Expliquez, puis corrigez.
\begin{bashcode}
while read -r ligne; do
    ssh cible "echo \$ligne"
done < liste.txt
\end{bashcode}
\begin{solution}
Le \texttt{ssh} de la première itération \emph{bavasse} dans le stdin hérité (la boucle
\texttt{read}) et consomme le reste du fichier. Correction : \texttt{-n}, qui redirige le
stdin de ssh vers \texttt{/dev/null} :
\begin{bashcode}
while read -r ligne; do
    ssh -n cible "echo \$ligne"
done < liste.txt
# ou, plus lisible :
while read -r ligne; do
    printf '%s\n' "$ligne" | ssh cible "echo stdin-reçu"
done < liste.txt
\end{bashcode}
\texttt{-n} est indispensable dans toute boucle qui appelle \texttt{ssh} avec un stdin
actif ; symétriquement, quand on \emph{veut} envoyer du stdin à la commande distante,
on utilise un pipe explicite.
\end{solution}
\question Transfert non interactif : utilisez \texttt{sftp} en mode batch pour déposer un
fichier sur la cible. Comparez avec \texttt{scp} (rappel : depuis OpenSSH 9.0,
\texttt{scp} utilise par défaut le protocole SFTP).
\begin{solution}
\begin{bashcode}
printf 'put rapport.pdf /tmp/\nbye\n' > depot.batch
sftp -b depot.batch cible
# ou une ligne :
sftp cible <<< $'put rapport.pdf /tmp/\nbye'
# scp équivalent (protocole SFTP depuis OpenSSH 9.0 ; -O pour revenir au legacy) :
scp rapport.pdf cible:/tmp/
\end{bashcode}
Pour des \emph{arborescences} nombreuses, \texttt{rsync -e ssh} reste l'outil de
référence (réincrémental, delta) ; \texttt{scp} est à réserver aux transferts simples.
\end{solution}
\question \textbf{Débogage} : connectez-vous avec \texttt{ssh -v bastion} et repérez dans
la sortie (i)~la clé de hôte proposée par le serveur, (ii)~la clé utilisée pour
l'authentification, (iii)~la méthode d'auth acceptée. À quoi sert \texttt{-vvv} ?
\begin{solution}
\begin{bashcode}
ssh -v bastion 2>&1 | grep -Ei "server host key|offering publickey|server accepts|authenticated"
# debug1: Server host key: ECDSA 256 SHA256:XXXX...
# debug1: Offering public key: ED25519 AAAA...
# debug1: Server accepts key: ...
# debug1: Authentication succeeded (publickey).
\end{bashcode}
\texttt{-v} : le niveau de debug standard, lisible. \texttt{-vvv} : tout le détail
protocole (chiffres, algorithmes négociés, chaque message) — utile pour un bug
d'interopérabilité ou un ticket à OpenSSH, pas pour le quotidien.
\end{solution}
\question \textbf{HOST KEY CHANGED} : changez l'IP NAT du bastion (ex. vers 2223) puis
reconnectez-vous sur l'ancienne IP. Que répond ssh ? Pourquoi ? Quelle est la procédure
légitime, et la différence entre \texttt{StrictHostKeyChecking=accept-new} et
\texttt{=no} ?
\begin{solution}
\begin{bashcode}
# après un changement d'IP/rebuild, l'ancienne IP renvoie une clé différente :
ssh -p 2222 root@127.0.0.1
# @@@@@@@@@@@@@@@@@@@@@@@@@@@@@@@@@@@@@@@@@@@@@@@@@@@@@@@@@@@
# @    WARNING: REMOTE HOST IDENTIFICATION HAS CHANGED!     @
# ... IT IS POSSIBLE THAT someone is mist...
# HOST KEY CHANGED!
\end{bashcode}
Le serveur présente une clé de hôte différente de celle enregistrée dans
\texttt{known\_hosts} : soit la VM a été reconstruite (cas légitime : nouveau
\texttt{sshd\_host\_key}), soit un MITM répond à la place du serveur. La procédure
légitime : vérifier l'empreinte \emph{out-of-band} (console physique, canal séparé),
puis :
\begin{bashcode}
ssh-keygen -R "127.0.0.1:2222"   # retire l'entrée obsolète de known_hosts
ssh -p 2222 root@127.0.0.1       # première connexion : acceptation consciente
\end{bashcode}
\texttt{accept-new} (OpenSSH 7.6+) : accepte et \emph{enregistre} les hôtes \emph{nuls},
mais refuse toujours un changement de clé existant — c'est le bon réglage par défaut
pour l'automatisation. \texttt{no} : accepte \emph{tout} silencieusement —
\textbf{jamais}, même « pour que le CI passe ».
\end{solution}
\question Un fichier \texttt{known\_hosts} peut être personnalisé avec
\texttt{KnownHostsFile}. Citez deux cas d'usage.
\begin{solution}
\begin{itemize}
  \item un service qui tourne sous un utilisateur \emph{dédié} (pas votre compte) avec
        ses propres clés : \texttt{KnownHostsFile /etc/ssh/known\_hosts\_service} ;
  \item des configs par \emph{projet} (une flotte de clients A vs de clients B) pour ne
        pas polluer le \texttt{known\_hosts} personnel.
\end{itemize}
\end{solution}
\question Votre connexion au campus coupe régulièrement (WiFi). Décrivez le rôle de
\texttt{ServerAliveInterval} / \texttt{ServerAliveCountMax}, puis montrez l'alternative
qui \emph{survit} aux pertes de réseau : \texttt{mosh}.
\begin{solution}
\begin{bashcode}
# côté config :
#   ServerAliveInterval 30   # envoie un keepalive chiffré toutes les 30 s
#   ServerAliveCountMax 3    # 3 keepalive sans réponse = connexion déclarée morte
\end{bashcode}
Ces keepalive \emph{détectent} les connexions mortes (NAT expiré, route disparue) et
provoquent une erreur nette plutôt qu'un hang. Mais SSH, en TCP, \emph{ne survit pas} au
changement d'adresse IP (l'uplet 4-tuple est cassé) : il faut reconnecter.
\texttt{mosh} (Mobile Shell), en UDP avec re-synchronisation d'état, \emph{survit} au
roaming et aux coupures courtes :
\begin{bashcode}
sudo apt install mosh
mosh bastion
# (test : décochez la carte NAT dans VirtualBox pendant 10 s, puis recochez-la)
# la session mosh est intouchable ; la session ssh classique est morte
\end{bashcode}
Pour du travail en salle de TP sur le WiFi, \texttt{mosh} + clés + config est le
combo le plus confortable ; pour de l'automatisation, on reste sur \texttt{ssh} avec
keepalive.
\end{solution}
\end{questions}

\section{Challenge : la config qui ne marche pas}

Votre collègue a laissé cette config. Il ne parvient pas à se connecter au serveur
interne \texttt{10.99.0.11} (alias \texttt{lab}) via le bastion. Il a en outre « perdu »
les permissions de sa clé en recopiant son \texttt{\textasciitilde/.ssh} sur un clé USB.

\begin{bashcode}
# ~/.ssh/config du collègue
Host lab
    HostName 10.99.0.11
    User root
    IdentityFile ~/.ssh/id_ed25519
    IdentitiesOnly yes
    ProxyCommand ssh -W %h:%p 127.0.0.1 -p 2222
    ControlMaster auto
    ControlPath ~/.ssh/ssh-multiplex-control-socket-for-the-connection-between-client-and-server-%r@%h-%p
    ControlPersist 10m
\end{bashcode}

Critère de réussite :
\begin{itemize}
  \item \texttt{ssh lab} fonctionne en une commande ;
  \item la \emph{seconde} \texttt{time ssh lab true} est inférieure à \textbf{50~ms} ;
  \item aucune clé n'est exposée sur le bastion (\texttt{ssh bastion ssh-add -l} n'affiche
    rien de plus que les clés du bastion).
\end{itemize}

\begin{questions}
\question Identifiez et corrigez les \textbf{trois} défauts. Pour chacun, donnez le
symptôme exact observé et la correction.
\begin{solution}
\textbf{1.~Permissions de la clé privée.} Recopié sur clé USB, le fichier
\texttt{id\_ed25519} est revenu en 644 :
\begin{bashcode}
# symptôme :
ssh lab
# @@@@@@@@@@@@@@@@@@@@@@@@@@@@@@@@@@@@@@@@@@@@@@@@@@@@@@@@@@@
# @         UNPROTECTED PRIVATE KEY FILE!          @
# Bad permissions.
chmod 600 ~/.ssh/id_ed25519   # correction
\end{bashcode}
\textbf{2.~\texttt{IdentitiesOnly yes} + seule la clé du bastion.} La config ne présente
qu'\texttt{id\_ed25519} (la clé du \emph{bastion}) aux deux sauts ; la cible n'a que
\texttt{id\_ed25519\_cible} dans son \texttt{authorized\_keys} :
\begin{bashcode}
# symptôme (après correction 1) :
ssh lab
# root@10.99.0.11: Permission denied (publickey).
# correction : présenter aussi la clé cible
IdentityFile ~/.ssh/id_ed25519_cible
\end{bashcode}
(Le \texttt{IdentitiesOnly yes} est ici ce qui \emph{empêche} l'agent de sauver la mise :
seules les clés listées sont présentées. On aurait pu le retirer, mais il est plus propre
de lister explicitement les clés de chaque saut.)

\textbf{3.~\texttt{ControlPath} trop long.} Le préfixe
\texttt{ssh-multiplex-control-socket-for-the-connection-between-client-and-server-}
pousse le chemin au-delà de 108~caractères :
\begin{bashcode}
# symptôme :
ssh lab
# .../id_ed25519:_SOCKET: Socket path too long
# correction : utiliser %C (hash borné)
ControlPath ~/.ssh/ctl-%C
\end{bashcode}
Config corrigée minimale :
\begin{bashcode}
Host lab
    HostName 10.99.0.11
    User root
    IdentityFile ~/.ssh/id_ed25519
    IdentityFile ~/.ssh/id_ed25519_cible
    ProxyCommand ssh -W %h:%p 127.0.0.1 -p 2222
    ControlMaster auto
    ControlPath ~/.ssh/ctl-%C
    ControlPersist 10m
\end{bashcode}
(Bonus : remplacer le \texttt{ProxyCommand} par \texttt{ProxyJump bastion} si un alias
\texttt{bastion} existe dans la config.)
\end{solution}
\end{questions}

\section{Cheatsheet}

\begin{center}
\small
\begin{tabular}{p{0.45\textwidth} p{0.5\textwidth}}
\hline
\textbf{Option} & \textbf{Usage} \\
\hline
\texttt{ssh-keygen -t ed25519 -C "comment"} & générer une paire de clés (algorithme recommandé) \\
\texttt{ssh-keygen -lf clé} & fingerprint (audit de \texttt{authorized\_keys}) \\
\texttt{ssh-keygen -y -f privée} & redérive la publique depuis la privée \\
\texttt{ssh-copy-id hôte} & déploie la clé publique \\
\texttt{ssh-keyscan -t ed25519 hôte} & pré-peuple \texttt{known\_hosts} (pas d'auth !) \\
\texttt{ssh-add [-t secondes]} & charge/décharge des clés dans l'agent \\
\texttt{ssh-add -l} & liste les clés de l'agent \\
\texttt{-i fichier} & choisit la clé présentée \\
\texttt{IdentitiesOnly yes} & n'utilise que les \texttt{IdentityFile} (pas l'agent) \\
\texttt{-A} / \texttt{ForwardAgent} & \textbf{piège} : expose l'agent sur le saut \\
\texttt{-J saut} / \texttt{ProxyJump} & connexion par saut (clé locale à chaque hop) \\
\texttt{-N} / \texttt{-fN} & pas de shell / tunnel en arrière-plan \\
\texttt{-L port:hôte:port} & tunnel local (hôte résolu \emph{du côté serveur}) \\
\texttt{-R port:hôte:port} & tunnel distant (écouteur \emph{sur} le serveur) \\
\texttt{-D port} & proxy SOCKS5 dynamique \\
\texttt{ssh -G hôte} & affiche la config résolue (debug) \\
\texttt{ssh -O check|cancel} & état / arrêt du maître multiplexé \\
\texttt{ControlMaster auto} & multiplexage sur socket UNIX \\
\texttt{ControlPath ~/.ssh/ctl-\%C} & chemin borné (\%C = hash) \\
\texttt{ControlPersist 10m} & survie du maître après déconnexion \\
\texttt{ServerAliveInterval 30} & keepalive : détecte les connexions mortes \\
\texttt{StrictHostKeyChecking accept-new} & accepte les nouveaux, refuse les changées \\
\texttt{ssh-keygen -R hôte} & retire une entrée de \texttt{known\_hosts} \\
\texttt{-v} / \texttt{-vvv} & debug (standard / protocole) \\
\texttt{mosh hôte} & shell qui survit au roaming (UDP) \\
\hline
\end{tabular}
\end{center}

\section{Pour aller plus loin}

\begin{itemize}
  \item \texttt{man ssh}, \texttt{man sshd\_config}, \texttt{man ssh-config} : la référence
        complète des options, y compris \texttt{Match} et les tokens
        \texttt{\%h}/\texttt{\%p}/\texttt{\%r}/\texttt{\%C} ;
  \item la spec du protocole : \emph{PROTOCOL} et \emph{PROTOCOL\.kex} dans la source
        OpenSSH ;
  \item \emph{Hardening OpenSSH} (guide de l'IETF) pour le durcissement
        \texttt{sshd\_config} côté serveur ;
  \item Mosh (\url{https://mosh.org/}) : le papier et l'implémentation ;
  \item l'option \texttt{PermitTunnel} de \texttt{sshd\_config} (tunnel TAP, pour du
        véritable VPN sur SSH) — à connaître, rarement à utiliser.
\end{itemize}

\end{document}
